\phantomsection
\label{fig:chu-han-corridors}
\begin{center}
\begin{tikzpicture}[x=.88cm,y=.88cm,font=\small]
  \fill[paper] (0,0) rectangle (12.3,6.55);
  \draw[source!70,line width=.8pt] (.45,5.32) .. controls (3.1,4.78) and
    (7.1,5.75) .. (11.85,4.88);
  \node[text=source,above,font=\scriptsize] at (8.4,5.2) {黄河方向（示意）};

  \fill[dynasty!25] (.85,.85) ellipse (1.3 and .86);
  \node[align=center] at (.85,.85) {关中\\汉方后方};

  \fill[timeline!26] (3.7,2.58) ellipse (1.13 and .78);
  \node[align=center] at (3.7,2.58) {荥阳—成皋\\对峙走廊};

  \fill[source!25] (5.55,4.25) ellipse (1.03 and .72);
  \node[align=center] at (5.55,4.25) {河东、赵地\\井陉方向};

  \fill[timeline!22] (8.45,4.42) ellipse (1.08 and .78);
  \node[align=center] at (8.45,4.42) {齐地\\潍水方向};

  \fill[dynasty!43] (9.15,1.65) ellipse (1.2 and .85);
  \node[align=center,text=white] at (9.15,1.65) {彭城\\西楚中心};

  \fill[source!18] (11.05,3.1) ellipse (.83 and .62);
  \node[align=center,font=\scriptsize] at (11.05,3.1) {垓下\\合围};

  \draw[-{Stealth[length=2mm]},very thick,color=dynasty]
    (2.04,1.22) .. controls (2.5,1.65) and (2.72,2.15) .. (2.88,2.36);
  \node[text=dynasty,below,font=\scriptsize] at (2.05,2.1) {关中补给};

  \draw[-{Stealth[length=2mm]},very thick,color=dynasty]
    (4.54,2.98) .. controls (4.92,3.5) and (5.08,3.62) .. (5.18,3.7);
  \node[text=dynasty,right,font=\scriptsize] at (4.38,3.85) {北进};

  \draw[-{Stealth[length=2mm]},very thick,color=dynasty]
    (6.52,4.34) -- (7.3,4.39);
  \node[text=dynasty,above,font=\scriptsize] at (6.9,4.4) {东进};

  \draw[-{Stealth[length=2mm]},thick,dashed,color=timeline]
    (8.8,3.65) .. controls (9.15,3.0) and (9.25,2.65) .. (9.24,2.56);
  \draw[-{Stealth[length=2mm]},thick,dashed,color=timeline]
    (10.05,2.32) .. controls (10.36,2.6) and (10.48,2.83) .. (10.55,2.93);
  \node[text=timeline,align=center,font=\scriptsize] at (10.47,1.45)
    {齐、梁等兵力\\转入合围};

  \draw[-{Stealth[length=2mm]},thick,dashed,color=source]
    (8.0,2.02) .. controls (6.6,2.13) and (5.45,2.38) .. (4.84,2.52);
  \node[text=source,above,font=\scriptsize] at (6.45,2.22) {楚军压力与争夺};

  \node[draw=source!75,fill=paper,rounded corners=2pt,align=center,
    inner sep=3pt,font=\scriptsize] at (2.0,5.74)
    {前205--前202年：\\后方、对峙走廊与北方战区相互牵动};
\end{tikzpicture}

\smallskip
\small\textit{图\quad 楚汉战争的几条战区联系，前205--前202年：据《史记·高祖本纪》《项羽本纪》《淮阴侯列传》及《汉书·高帝纪》《韩信传》，提示关中后方、荥阳—成皋对峙、赵齐战区与垓下合围之间的资源联系。河流、节点和箭线均为约略示意；不复原军队实际行军轨迹、战场位置、补给路线或任何精确政治边界。}
\end{center}
